\documentclass[10pt,a4paper]{amsart}
\usepackage{amsmath,amssymb,amsthm,graphicx}
\usepackage{verbatim}
\newenvironment{reference}{\comment}{\endcomment}
%\newenvironment{reference}{}{}
\newenvironment{slogan}{\comment}{\endcomment}
\newenvironment{history}{\comment}{\endcomment}

% For commutative diagrams we use Xy-pic
\usepackage[all]{xy}

% We use 2cell for 2-commutative diagrams.
\xyoption{2cell}
% Theorem environments.
%
\theoremstyle{plain}
\newtheorem{theorem}[subsection]{Theorem}
\newtheorem{proposition}[subsection]{Proposition}
\newtheorem{lemma}[subsection]{Lemma}

\theoremstyle{definition}
\newtheorem{definition}[subsection]{Definition}
\newtheorem{example}[subsection]{Example}
\newtheorem{exercise}[subsection]{Exercise}
\newtheorem{situation}[subsection]{Situation}

\theoremstyle{remark}
\newtheorem{remark}[subsection]{Remark}
\newtheorem{remarks}[subsection]{Remarks}

\numberwithin{equation}{subsection}

% Macros
%
\def\lim{\mathop{\mathrm{lim}}\nolimits}
\def\colim{\mathop{\mathrm{colim}}\nolimits}
\def\Spec{\mathop{\mathrm{Spec}}}
\def\Hom{\mathop{\mathrm{Hom}}\nolimits}
\def\Ext{\mathop{\mathrm{Ext}}\nolimits}
\def\SheafHom{\mathop{\mathcal{H}\!\mathit{om}}\nolimits}
\def\SheafExt{\mathop{\mathcal{E}\!\mathit{xt}}\nolimits}
\def\Sch{\mathit{Sch}}
\def\Mor{\mathop{\mathrm{Mor}}\nolimits}
\def\Ob{\mathop{\mathrm{Ob}}\nolimits}
\def\Sh{\mathop{\mathit{Sh}}\nolimits}
\def\NL{\mathop{N\!L}\nolimits}
\def\CH{\mathop{\mathrm{CH}}\nolimits}
\def\proetale{{pro\text{-}\acute{e}tale}}
\def\etale{{\acute{e}tale}}
\def\QCoh{\mathit{QCoh}}
\def\Ker{\mathop{\mathrm{Ker}}}
\def\Im{\mathop{\mathrm{Im}}}
\def\Coker{\mathop{\mathrm{Coker}}}
\def\Coim{\mathop{\mathrm{Coim}}}

% Boxtimes
%
\DeclareMathSymbol{\boxtimes}{\mathbin}{AMSa}{"02}

%
% Macros for moduli stacks/spaces
%
\def\QCohstack{\mathcal{QC}\!\mathit{oh}}
\def\Cohstack{\mathcal{C}\!\mathit{oh}}
\def\Spacesstack{\mathcal{S}\!\mathit{paces}}
\def\Quotfunctor{\mathrm{Quot}}
\def\Hilbfunctor{\mathrm{Hilb}}
\def\Curvesstack{\mathcal{C}\!\mathit{urves}}
\def\Polarizedstack{\mathcal{P}\!\mathit{olarized}}
\def\Complexesstack{\mathcal{C}\!\mathit{omplexes}}
% \Pic is the operator that assigns to X its picard group, usage \Pic(X)
% \Picardstack_{X/B} denotes the Picard stack of X over B
% \Picardfunctor_{X/B} denotes the Picard functor of X over B
\def\Pic{\mathop{\mathrm{Pic}}\nolimits}
\def\Picardstack{\mathcal{P}\!\mathit{ic}}
\def\Picardfunctor{\mathrm{Pic}}
\def\Deformationcategory{\mathcal{D}\!\mathit{ef}}
\usepackage{cleveref}
\crefname{originallemma}{original lemma}{original lemmas}
\Crefname{originallemma}{Original lemma}{Original lemmas}
\makeatletter
\newlabel{lemma-directed}{{00X3}{source}{}{sourcepointer-1}{}}
\newlabel{lemma-directed@cref}{{[originallemma][0][]00X3}{[1][0][]source}{}{}{}}
\makeatother
\crefname{originallemma}{original lemma}{original lemmas}
\Crefname{originallemma}{Original lemma}{Original lemmas}
\makeatletter
\newlabel{categories-lemma-finite-limits-exist}{{002O}{source}{}{sourcepointer-2}{}}
\newlabel{categories-lemma-finite-limits-exist@cref}{{[originallemma][0][]002O}{[1][0][]source}{}{}{}}
\makeatother
\crefname{originallemma}{original lemma}{original lemmas}
\Crefname{originallemma}{Original lemma}{Original lemmas}
\makeatletter
\newlabel{lemma-neighbourhoods-cofiltered}{{0F4E}{source}{}{sourcepointer-3}{}}
\newlabel{lemma-neighbourhoods-cofiltered@cref}{{[originallemma][0][]0F4E}{[1][0][]source}{}{}{}}
\makeatother
\title{Conditions ensuring that the category of neighbourhoods is directed}
\author{The Stacks Project Authors}
\date{stacks-project:sites.tex:local-snapshot-20261012; GNU FDL 1.2 or later}
\begin{document}
\maketitle
\noindent\textit{Editorial note.} All delivered mathematical text is the author's own, copied from the pinned Stacks Project chapter file sites-cohomology.tex; no statement or proof step was rewritten, reordered or removed.\par\medskip
\noindent\textit{Editorial note.} The proof is the author's complete argument; every step the author writes out is inside the counted span and no closing sentence is a placeholder.\par\medskip
\noindent\textit{Editorial note.} Every labelled reference occurring inside the counted spans is carried in assembly.external\_references as a non-load-bearing navigation pointer (same-chapter ones at their definition site, cross-chapter ones at the invocation inside the counted span), so no load-bearing mathematics is replaced by a stub. No bibliography citation occurs inside the counted spans.\par\medskip
\noindent\textit{Editorial note.} The checked chapter list for same-name collisions is the set of stacks chapters staged in this repository: chow, cohomology, divisors, etale-cohomology, flat, more-algebra, more-morphisms, morphisms, spaces-more-morphisms, varieties together with this card's chapter sites-cohomology.\par\medskip
\noindent\textit{Editorial note.} The preamble is the chapter's own retained minimal preamble (sources/stacks-sites-cohomology/retained-preamble.tex), a verbatim excerpt of the native preamble: the Xy-pic setup and the native theorem declarations with the mathematical macros.\par\medskip
\noindent\textit{Allowed standard prerequisites (editorial).}\begin{itemize}
\item Existence of finite limits. Categories, Lemma categories-lemma-finite-limits-exist. Supplies the construction of equalizers from fibre products.
\item Cofiltered neighbourhoods imply left exactness of the stalk functor. Same chapter, Lemma lemma-neighbourhoods-cofiltered. Supplies the concluding clause of the statement.
\end{itemize}\medskip
\ifdefined\hypertarget\hypertarget{sourcepointer-1}{}\fi
\noindent\textit{Source reference.} Original lemma 00X3: sources/stacks-sites/pinned-original-source.tex, source lines 8501--8501. Navigation pointer invoked at proof line 8501; the counted span contains the pointer only and the cited result is not reproduced inside it, so the pointer is registered as a non-load-bearing navigation stub in the established form for Stacks material. No bibliography citation is involved.\par\medskip
\ifdefined\hypertarget\hypertarget{sourcepointer-2}{}\fi
\noindent\textit{Source reference.} Original lemma 002O: sources/stacks-sites/pinned-original-source.tex, source lines 8503--8503. Navigation pointer invoked at proof line 8503; the counted span contains the pointer only and the cited result is not reproduced inside it, so the pointer is registered as a non-load-bearing navigation stub in the established form for Stacks material. No bibliography citation is involved.\par\medskip
\ifdefined\hypertarget\hypertarget{sourcepointer-3}{}\fi
\noindent\textit{Source reference.} Original lemma 0F4E: sources/stacks-sites/pinned-original-source.tex, source lines 8522--8522. Navigation pointer invoked at proof line 8522; the counted span contains the pointer only and the cited result is not reproduced inside it, so the pointer is registered as a non-load-bearing navigation stub in the established form for Stacks material. No bibliography citation is involved.\par\medskip
\begin{lemma}
\label{lemma-neighbourhoods-directed}
Let $\mathcal{C}$ be a site. Assume that $\mathcal{C}$ has
a final object $X$ and fibred products.
Let $p = u : \mathcal{C} \to \textit{Sets}$ be a functor such that
\begin{enumerate}
\item $u(X)$ is a singleton set, and
\item for every pair of morphisms $U \to W$ and $V \to W$ with
the same target the map
$u(U \times_W V) \to u(U) \times_{u(W)} u(V)$ is bijective.
\end{enumerate}
Then the the category of neighbourhoods of $p$ is cofiltered
and consequently the stalk functor $\Sh(\mathcal{C}) \to \textit{Sets}$,
$\mathcal{F} \to \mathcal{F}_p$ commutes with finite limits.
\end{lemma}

\begin{proof}
Please note the analogy with Lemma \ref{lemma-directed}.
The assumptions on $\mathcal{C}$ imply that $\mathcal{C}$ has finite limits.
See Categories, Lemma \ref{categories-lemma-finite-limits-exist}.
Assumption (1) implies that the category of neighbourhoods
is nonempty. Suppose $(U, x)$ and $(V, y)$ are neighbourhoods.
Then
$u(U \times V) = u(U \times_X V) =
u(U) \times_{u(X)} u(V) = u(U) \times u(V)$ by (2).
Hence there exists a neighbourhood $(U \times_X V, z)$ mapping
to both $(U, x)$ and $(V, y)$.
Let $a, b : (V, y) \to (U, x)$ be two morphisms
in the category of neighbourhoods. Let $W$ be the equalizer of
$a, b : V \to U$. As in the proof of
Categories, Lemma \ref{categories-lemma-finite-limits-exist}
we may write $W$ in terms of fibre products:
$$
W = (V \times_{a, U, b} V) \times_{(pr_1, pr_2), V \times V, \Delta} V
$$
The bijectivity in (2) guarantees there exists an element $z \in u(W)$
which maps to $((y, y), y)$.
Then $(W, z) \to (V, y)$ equalizes $a, b$ as desired.
The ``consequently'' clause is Lemma \ref{lemma-neighbourhoods-cofiltered}.
\end{proof}

\end{document}
